\section{Sicherheit}
\label{sec:safety}

\subsection{Bestimmungsgemässe Verwendung}
Der \ST{} ist ein tragbares, batteriebetriebenes Prüfgerät für die elektrische Abnahme von
\ac{FS}-Fahrzeugen. Er darf verwendet werden, um
\begin{itemize}
    \item einen Prüfling mit einer Gleichspannung (\qtyrange{0}{60}{\volt}) oder einem Gleichstrom (\qtyrange{0}{10}{\ampere}) zu versorgen,
    \item kleine Widerstände in Vierleiterschaltung und Widerstände in Zweileiterschaltung zu messen,
    \item Isolationswiderstände mit Prüfspannungen von \qtylist{125;250;500;1000}{\volt} zu messen,
    \item Gleichspannungen und Gleichströme zu messen.
\end{itemize}
Jede andere Verwendung gilt als nicht bestimmungsgemäss. Das Gerät darf nur von Personen
verwendet werden, die für Arbeiten an elektrischen Anlagen qualifiziert sind und diese
Anleitung gelesen haben.
\todo[inline]{Erforderliche Qualifikation festlegen (z.\,B. FS ESO / Scrutineer, elektrotechnisch unterwiesene Person)}

\subsection{Allgemeine Sicherheitshinweise}

\begin{danger}
    In der Betriebsart \disp{Isolation Test} liegen an den Prüfbuchsen bis zu \qty{1000}{\volt} DC an.
    Prüfleitungen, Buchsen und Prüfling während der Prüfung nicht berühren.
    Der Prüfling kann nach der Prüfung geladen bleiben.
    \todo[inline]{Bestätigen: wird der Prüfling nach der Isolationsprüfung aktiv entladen (Q18)? Wie lange warten?}
\end{danger}

\begin{warning}
    \begin{itemize}
        \item Gerät nur mit unbeschädigtem Gehäuse, unbeschädigten Buchsen und Prüfleitungen verwenden.
        \item Nur Prüfleitungen mit berührungsgeschützten \qty{4}{\milli\meter}-Steckern verwenden, die für die Spannung der Messung zugelassen sind.
        \item Messkategorie und maximale Eingangsspannung nicht überschreiten (\autoref{sec:specs}).
        \item In den Quellen-Betriebsarten, den Widerstandsmessungen, \disp{Ampmeter} und
              \disp{Isolation Test} keine Fremdspannung an die Buchsen anlegen.
        \item Gerät nicht in der Nähe brennbarer Gase oder Dämpfe verwenden.
    \end{itemize}
    \todo[inline]{Messkategorie (Anforderung: CAT III 600 V) noch nicht verifiziert / zertifiziert}
\end{warning}

\begin{caution}
    \begin{itemize}
        \item Gerät nur wie in \autoref{sec:charging} beschrieben laden.
        \item Das Gerät enthält eine \ac{LiFePO4}-Batterie. Nicht öffnen, kurzschliessen oder erhitzen.
        \item Betriebstemperatur: \tbc{\qtyrange{0}{50}{\celsius}}. Schutz gegen Regen: \tbc{IP54}.
    \end{itemize}
\end{caution}

\subsection{Schutzfunktionen}
Das Gerät enthält mehrere unabhängige Schutzfunktionen. Sie schützen das Gerät; sie ersetzen
nicht die Sicherheitsregeln für Arbeiten an elektrischen Anlagen.
\begin{itemize}
    \item \ac{BMS}: schützt die Batterie vor Über-/Unterspannung, Überstrom und Übertemperatur.
    \item Überspannungsschutz: wird eine zu hohe Spannung am Ausgang erkannt, wird der Ausgang
          kurzgeschlossen und das Gerät gesperrt, bis es aus- und wieder eingeschaltet wird
          (\autoref{sec:messages}). Bei einer externen Quelle kann dabei die interne Sicherung auslösen.
    \item Überstromschutz des \ac{HV}-Ausgangs.
    \item Spannungs- und Stromgrenzen je Betriebsart, Temperaturüberwachung.
    \item Automatische Abschaltung nach \qty{10}{\minute} ohne Bedienung.
\end{itemize}
\todo[inline]{Not-Aus: Hardware-Abschaltung über ESC beschreiben (wie lange halten?)}
